\chapter{The Ordinary World}
\label{chap:pledge-ch1}

\section{Compositionality}
\label{sec:pledge-i}

There are so many ways into this topic that i find myself suffering a bit of option paralysis. What do computer and telecommunications networks, quantum mechanics and gravitation, scientific reductionism, and fractals have to do with each other? One loose knit community of computer scientists, category theorists, logicians, and type theorists  calls this phenomenon ``compositionality.'' Conceived in this way it's about a radical commitment to reasoning about and constructing larger, more complex systems by assembling or \emph{composing} smaller, simpler systems. For example, building a computer network by assembling smaller subnetworks, or building a cell out of molecules, or building a company out of divisions.

Indeed, what falls under the rubric ``scientific reductionism'' is actually just compositionality in disguise. The ``software stack'' of Western science is that ecosystems are composed of populations of species. A population of a given species is composed of individual organisms. An organism is composed of cells. Cells are composed of molecules. Molecules are composed of atoms. Atoms are composed of protons, neutrons, and electrons. These are composed of quarks. (As Nima Arkani-Hamed points out, what comes after this is effectively hidden behind the event horizon of a black hole because --- roughly speaking --- the amount of energy necessary to make observations at the Planck scale creates a black hole.)

Suffice it to say, somewhere in our psychology, or in our bones we get this idea. And yet we don't get this idea. For example, one way to see this is in the failure of quantum mechanics to talk to general relativity. In pop science versions quantum mechanics is about ``small scale'' stuff while general relativity is about ``large scale'' stuff. If that were so, then shouldn't the large scale stuff emerge from the right assemblies of the small scale stuff? 

In point of fact, in their original conceptions neither theory is what we call compositional. Einstein's equations are formulated for the entire universe, not local regions. To bridge between Einstein's theory and Schroedinger's you also need a wave equation for the entire universe.  Neither of these is available to humanity's purview. Humanity needs to piece together bigger pictures from what we can see locally. More on this later.

Another way to understand how we humans both get and don't get the idea of compositionality is in how we coordinate with each other. In some very real sense humanity's superpower is coordination. A single human is no match for a whale. But a group of skilled fisherfolk working in coordination can make relatively swift work of dispatching such a large sea creature. Indeed, these majestic creatures are endangered at the species level because of humanity's ability to coordinate. Somehow a group of humans working together can assemble a ``super human.''

Yet humans have a hard time coordinating at levels beyond the tribal. We have a hard time assembling a group of ``super humans'' (like a tribe) into a ``super super human.'' The current global order was achieved through vast bloodshed and is far from stable. Tribal fault lines reassert themselves over and over again into political  discourse, and yes, even academic discourse. In archeology, for example, going against the ``Clovis first'' clique was death to one's career for decades, and is only now crumbling under the evidence. Examples from the history of STEM of this human foible are a dime a dozen. No part of the history of human institutions is free of this phenomenon.

This fractured understanding of compositionality --- where we both get it and we don't --- also informs the way we build our tools. Computers and computer networks provide a fantastic illustration. There was a time when a computer was a lone entity. At first these lone entities were behemoths occupying entire warehouses. It didn't take very long, relatively speaking, for them to shrink to the size where they could fit on top of a desk. Yet, still they were lone operators. They didn't talk to each other.

Indeed the time from the first mainframe computers to the desktop was about the same as it was from the desktop to the modern Internet integrating mobile telecommunications. The evolution from ENIAC-class computers to the first Apple desktop computers took about 30 years. ENIAC was completed in 1945, while the first Apple desktop, the Apple I, was launched in 1976. As for the transition from the first Apple desktops to the modern Internet integrating mobile telecommunications, this period spans roughly 20 years. The Apple Macintosh (a significant early desktop) was introduced in 1984, and mobile telecommunications began integrating with the Internet in the early 2000s. So, ENIAC to Apple I took 30 years, and Apple I to Internet with mobile telecommunications took about 30 years. 

Yet, if we look inside these devices what do we see? \emph{A network of components}. And if we look at more modern processors like the Groq or Cerebras chips, what do we see? It begins to resemble Internet on a chip. Even further, field programmable gate arrays (FPGAs) model the kind of mobility and fluidity we see at the Internet level, where both individual devices and entire generations of devices are coming and going. And at the other end of the telescope, in our era where we are seeing the emergence of AI, the understanding guiding both the abstraction (neural networks) and the hardware that implements the abstraction is that \emph{the network is the computer}.

When we put these two ideas together, just like adjusting our side view and rear view mirrors and occasionally swiveling our heads, we can see into our blind spot. The network is the computer is the network is the computer is \ldots{} . This ``recursion'' loops us back to the formulation of compositionality in the first paragraph of this section. Are the systems we build by composing a collection of component systems actually simpler? They might be smaller, but are they actually simpler?

This question is what connects our investigation to fractals. Fractals are instances of systems that enjoy a kind of self-similarity. When you compose copies of these kinds of systems you get a sort of transformed --- yet somehow the same --- version of any one of the copies that you used to assemble the system. There's a lot more to explore here and we will in the upcoming chapters. But, i want to whet your appetite for the material by pointing out something that i believe is critical. Time. 

When we get a proper account of compositionality, a different picture of time and how it relates to system dynamics emerges. Time emerges, in part, as depth in the exploration of the composite structure. Not just depth in the structure, but depth in the activity of disassembling and reassembling the composite structure. Not to put too fine a point on it, but this is a radical shift in perspective.

Let's consider this in contrast to modern physics. In modern physics time is exogenous to the system. We set up some experiment and we use devices like clocks, outside of the experiment, to measure the passage of time. Part of the real challenge of the maths of general relativity is what to do about different experiments in different labs using different clocks. How are we to understand how they are synchronized? Does there have to be a global clock ticking down the seconds in the life of the universe? The so-called ``fabric of spacetime'' is really a representation of an exogenous framework of measurement devices, rods and clocks as Wheeler called them. But rods and clocks are physical devices and so need to be \emph{inside} the model, not outside it.

When we get the story about compositionality on better footing, when we see into our blind spot, what happens is that time is endogenous to the model of the system and its dynamics. It emerges naturally from the theory of causation that the system dynamics is exploring. But now i'm getting ahead of myself. Hopefully, this discussion motivates you Dear Reader to go with me to see how this works out in detail.

\section{Agency}
\label{sec:pledge-ii}

Underlying the naive story of composing ``super humans'' from collections of humans is a more subtle story about how agents compose. The central question is: ``Is agency really compositional?'' The answer in this book is a resounding ``Yes!'' But, it's going to take us a while to get there. One way to understand the proposal is to look at the phenomenon at three different scales of human experience: human activity at the global scale involving multiple countries; human activity at the scale of a single country; and human activity at the level of an individual human. What i'm going to argue is that there is simultaneously both more and less agency than we typically attribute to the constituents we identify as ``acting'' in those circumstances.

For the first example, we'll consider a theater of active conflict during World War II. For the second, we'll consider US Congressional deliberation. For the third, we'll look at how an individual human's decisions and decision making might change over the course of a day. What we see over and over again is that there is an ascription of agency to an entity that is more narrative than substance. Meanwhile what is tangible, observable, measurable is a distribution of resources.

History has it that in the D-Day operations the allied forces landing at Utah, Omaha, Gold, Juno, and Sword took different approaches related to their different resource constraints. What's crucial here is the status of the agent understood in our narratives of the history of WWII as the ``allied forces.'' We can say things like the allied forces won at Normandy, but the reality was that the US and British and Canadian forces, despite being rolled up under Eisenhower for that campaign, were faced with different conditions. The British though as brave as anyone else in the field were typically more cautious about the deployment of resources because they were more resource constrained. The US forces were more aggressive about the deployment of resources because they had a much more massive economic engine behind their fortification. Thus, the agency of the allied forces is largely post facto. We see it after key decisions about resources have been made and after key actions have been executed by the entire collective.

We find the same is true in the US Congress. Neither Republicans nor Democrats are necessarily unified agents. That narrative of the agency of the Republican party or the Democratic party arises after key events like a vote. Indeed the agency of the House as a whole is a post facto narrative. We say the US House approved or denied a bill after a vote. The vote represents a commitment to the deployment of certain resources under the control of the individual representatives.

None of this is particularly surprising. We navigate this understanding of agency all day everyday. The key question is why should it be any different for an individual human? Is there really an agent there? The narrative i want to put forward is that much like there are physical devices making up your computer or your phone, like a camera, and a microphone, and a display, etc and these devices are paired up with software representations ``device drivers'' and ``event handlers'', there are physical subsystems of a human body: muscles, tendons, and bones (actuators and servo-motors) and eyes and ears and nose and skin and tongue (sensors) paired up with ``software'' level device drivers and event handlers. These are the resources that get committed. And there are software and hardware pairs connected with longer term commitment of these resources that we often refer to as drives, such as reproductive drives, or metabolic drives. 

Just like in the European theater during WWII or in the US House of Representatives there are lots of ``agents'' in the human mind/body that have their own agendas. They make coalitions and these coalitions battle it out for access to the resources. When a particular coalition ``wins'' the human ``agent'' lurches in the direction set by the agenda of that coalition. For example reproductive and metabolic cabals win in two humans and Alice and Bob agree to a date at a restaurant. Just as with the allied forces or the Democrats or the House, the agency of Alice or Bob is post facto. It comes after the commitment of resources determined by the coalition of both inner and external ``agents.'' And, after the sun goes down and the process of digesting the meal begins, different coalitions of inner and outer agents win the allocation of the resources of muscles, skin, bones, senses. So Alice and Bob lurch in a different direction.

What appears to be a single agent, like Alice, is really only the history or trace of a series of commitments of resources. Alice went to Oberlin, then she took a year off and traveled in South America, then she went to graduate school at the University of Chicago where she met Bob. Etc. One thing that complicates this rather simple story about agency is that there is an obvious evolutionary advantage to agents like Alice or Bob building up a software model of both Alice and Bob. The key question is whether these models are actually in the driver's seat, actually providing some kind of central control, or whether they are on somewhat different footing.

Consider, even Eisenhower is not directly in control of the boots on the ground at Omaha beach. There is a complex arrangement of consent across multiple levels of ``agency'' that gets narrated as ``Eisenhower commanding the allied forces landing at Omaha beach.'' Likewise, we can consider a range of possibilities for the simulation of Alice running on Alice's wetware. For example, we don't think of a database as having much agency, rather we users of the database query it to reveal certain relationships in the data. Similarly, a simulation doesn't have much agency. When we code up a simulation of the solar system or the weather patterns we users use it to query possible futures of these systems. Alice's internal model of either Alice or Bob might be just such a system, a database or a simulation, the processing on which might be committed by a winning coalition of agents running on Alice's brain.

Note that in this scenario none of the coalitions of agents running on Alice's wetware necessarily have to follow the predictions made by the Alice simulation. Perhaps each one of us has had an experience something like this: i knew better, but i did it anyway. Alice's ``i'' in situations like this could well be a complex combination of a simulation of Alice (and indeed Bob, and/or Charlie, and \ldots{}) as well as a record of the commitments of resources and actions taken. Furthermore, even if a majority of coalitions of agents running on Alice's wetware determine that the Alice simulation is highly accurate and should be deferred to, that still doesn't mean that the Alice simulation is ``in charge.'' In fact, in that scenario the Alice simulation isn't even in charge in a manner analogous to the sense in which the Speaker of the House is in charge of the House, or Eisenhower was in charge of the allied forces.

Before going deeper into this, however there is another crucial question just begging to be brought into the conversation. The Alice simulation does not necessarily have to be wired up with the other agents running on Alice's wetware in such a way that Alice has an experience of what it's like to be Alice. That's another layer of complexity and architecture that we will tackle in the upcoming sections and show that this is closely related, but not the same as the Alice simulation having an agency on par with Alice's reproductive or metabolic drives.

For now, i just want to leave you Dear Reader, with the picture of the wetware of a human (or indeed any creature on planet Earth) as the potential host of a variety of software agents, independent software agents. These software agents are not necessarily aligned. Rather, they compete for --- race for --- the resources available in the wetware, and are potentially also resources themselves for other software agents in the mix. The winners of these competitions --- these races --- determine a committed deployment of the actual resources. What we think of as agency --- at any level of this hierarchy --- is really a post facto narrative. It comes as a result of a history or trace of commitments of resources. From this picture, and its variants, we can build up a compositional account of agency.

\section{The Bootstrapping Problem}
\label{sec:pledge-iii}

The mathematician and computer scientist in me is itching to lay out the technical framework that makes the conversation many times more precise. Yet, another part of me is well aware that as soon as the narrative becomes technical i will lose a significant percentage of the readership. A coalition of agents running on my wetware is therefore making a gambit: defer the deployment of the technical apparatus; find ways to develop more of a relationship and more trust with the readers. So that's what ``i'm'' going to do.

One of the ways to develop trust is to show some vulnerability. Show something personal. Arguably, it's less personal and less vulnerable if i have already revealed an ulterior motive. Perhaps, though, that is also part of the gambit these agents in ``me'' are making. Being radically honest about ``my'' motives is also a way of being vulnerable. Either way, i want to mention something about my own process over the last few years. After the failure of my last marriage i decided to take stock of myself and my relationships and noticed something striking. 

Without exception, the people with whom i had intimate relationships were very much at the far perimeter of my life. Certainly not present in my day to day activities. Meanwhile the people with whom i have been friends, and not intimate partners, are still my friends. Still very much closer to the center of my present moment. Now, i have friends who successfully maintain and manage meaningful and fulfilling relationships with their exes. On the other hand, i appear to have failed twice: both in maintaining and developing the initial form of the relationship and in maintaining subsequent forms.

Having noticed this about myself i decided to investigate this phenomenon much more deeply and eschew intimate involvement until i had a deeper understanding of these dynamics in myself. One of the observations that came out of this investigation is that i realized that no relationship between two humans is just between those two humans. Each human, on average, comes equipped with \emph{two} copies of the really important humans in their lives. For example, each human has a mother and a father. Statistically speaking, for an adult human --- let's stick with Alice as a name --- younger than middle age, that human's progenitors are still alive. Moreover, if her upbringing involved both parents over a long enough time, then Alice has developed a simulation --- at least a narrative, but often other kinds of modeling, as well --- of both parents running in her wetware. This is also true of siblings, grandparents, teachers, friends, etc.

When Alice meets Bob they bring to the relationship their respective networks of actual family and friends, and their simulations of those networks running in their respective wetware. If Alice and Bob become significant to each other, if they begin to share their narratives about their parents, their siblings, their friends, their coworkers, as well as introduce each other to their actual networks, then they begin to run simulations of each other's networks, as well as having to participate in this much wider network of relationships. Alice runs not only simulations of her community of significants, but also Bob's network. And likewise for Bob. This is astonishingly hard to navigate.

For example, Alice may notice that Bob's narrative or simulation of his mother doesn't match Alice's experience of Bob's mother. Conversely, Alice may interpret her experience of Bob's mother through Bob's account of his simulation of his mother --- from which her simulation of Bob's mother was derived before she met Bob's parents. The opportunities for friction and conflict and misunderstanding abound. And, yet humans navigate these complexities with surprising deftness.

Part of this has to do with careful resource management. There's an exponential explosion of potential reflective modeling. Alice models Bob modeling his mother and compares it to her own model of Bob's mother. This has to be curtailed rather aggressively even if Alice has a large cranium and excellent compression techniques. One of the ways to do this is to aggressively avoid copying wherever possible.

What does this mean? There are two phenomena that interact that certainly generate copies. When simulations involve non-determinism, and when the simulations are used to look relatively far into the future. For example, in Alice's experience, when faced with certain situations Bob's mother may either do A or do B. And Alice's simulation doesn't favor either outcome. Indeed, the outcomes might have to do with some external non-determinism, such as the weather. As Alice looks into the future of her simulation of Bob's mother she has to carry around both the A outcome and the B outcome for Bob's mother. Unless Alice is careful, she is potentially making copies of a significant portion of the state she maintains about Bob's mother (her hair color, her educational history, etc) for both the A circumstances and the B circumstances. The situation gets worse the farther Alice looks into a future involving non-determinism. Further exploration involves maintaining more branches and potentially more copies.

It turns out there is a canonical way to compress this information so that all things being equal Alice's simulations are optimally organized. Taking a step back, the picture that emerges is one of interacting networks, rather than single agents. This is consistent with the theme we explored in the previous section. Agency is a post facto narrative that emerges as a history of deployments of resources. When we peal back the narrative of agency we find that it is rooted in the interaction of networks of agents, which are in turn rooted in the interaction of networks of agents, etc. This recursive picture is closely related to what we discussed about fractals and self-similarity. In upcoming chapters we will spell this out in glorious, mathematically precise detail.

But, i'm imaging that you, Dear Reader, are wondering: has this understanding helped me in my process? i cannot honestly say. What i can say is that the process has led me to one of the most significant and satisfying relationships of my life. One that may have a staying power bigger than my tendencies to fall apart. More on this later.
